% Aarons conclusion
\section{Conclusion}
This paper applied the complexity-adjusted benchmark of \textcite{ludwig2026local} to the
finite-sample comparison of Vector Autoregressions and Local Projections, replacing the
conventional equal-lag contest with one that holds effective complexity constant. Under correct
specification a fixed $LP(4)$ already behaves like a highly parameterized $VAR(q \approx 15)$, so
the familiar point-estimation edge of low-order VARs is largely an artifact of parsimony rather
than an intrinsic advantage of the iterated method. Functional misspecification, introduced
through a nonlinear lag-quadratic process, leaves this reading intact. Because the centered
quadratic is approximately orthogonal to the linear lag space, the estimand stays nearly invariant
to the nonlinearity, and asymptotically both $LP(4)$ and the adequately lagged VARs collapse onto
the same best linear approximation, so the population equivalence of \textcite{plagborg2021local}
survives without an asymptotic bias floor. In finite samples the $VAR(4)$ keeps its mean squared
error dominance, since the direct horizon-by-horizon regressions of the local projection carry a
larger variance.

The advantage that is absolute is inferential. The unmodeled quadratic injects conditional
heteroskedasticity into the residuals, which the homoskedastic delta-method intervals of the VAR
cannot capture, so they understate the sampling dispersion and under-cover ever more severely as
persistence and nonlinearity rise, while the heteroskedasticity-robust intervals of $LP(4)$ stay
near nominal. This is the first of our designs in which the robustness of local projections becomes
a tangible advantage, though it is confined to inference. The broader lesson is that estimand
equivalence in population does not imply inferential equivalence in finite samples, and the choice
between the two methods is objective-dependent, weighing the efficiency of the VAR point estimate
against the interval validity of the direct projection.

Three limitations bound these conclusions. First, the full VAR sweep up to $VAR(23)$ makes smaller 
samples infeasible, since the high-order systems turn rank-deficient and unstable,
leaving the small-sample region, where the bias-variance trade-off is sharpest and an LP advantage
most plausible, unexplored. Second, the nonlinear lag-quadratic destabilizes the simulated paths at
high persistence, so we restricted persistence to $\rho \in \{0.5, 0.7\}$ and left the
near-unit-root region unexamined, precisely where the interaction of high persistence and
functional misspecification would most strain inference. Third, the VAR coverage failure reflects
the default homoskedastic delta-method standard error rather than the iterated architecture itself,
and equipping the VAR sweep with robust or wild-bootstrap standard errors would likely narrow the
inferential gap.

Future work should extend the complexity-adjusted frontier to nonlinear settings that permit high
persistence without explosive paths, so the near-unit-root region can be revisited, and should
formalize a complexity adjustment for bootstrap-based inference. The latter would deliver a cleaner
horse race, isolating whether the robust coverage of local projections is an intrinsic virtue of
direct multi-step forecasting or merely a mechanical benefit of robust variance estimation.